\documentclass[11pt]{article}
\usepackage[margin=1in]{geometry}

% Core packages
\usepackage{amsmath,amssymb,amsthm}
\usepackage{hyperref}

% Paragraphs
\setlength{\parindent}{0pt}
\setlength{\parskip}{1\baselineskip}

\title{Expectation Maximization (EM)}
\author{Benjamin Yeh}
\date{2026-04-21}

\theoremstyle{definition}
\newtheorem*{theorem}{Theorem}
\newtheorem*{definition}{Definition}
\newtheorem*{remark}{Remark}

\DeclareMathOperator*{\argmax}{arg\,max}
\DeclareMathOperator*{\argmin}{arg\,min}

\begin{document}
\maketitle

References
\begin{itemize}
    \item Wikipedia: \url{https://en.wikipedia.org/wiki/Expectation%E2%80%93maximization_algorithm}
    \item Stanford CS229 notes: \url{https://cs229.stanford.edu/main_notes.pdf}
    \item Caltech BE/Bi 103b: \url{https://bebi103b.github.io/lessons/optimization/em_algorithm.html}
\end{itemize}

\section{Background}

\subsection{Theorems and Definitions}

\begin{theorem}[Law of the unconscious statistician]
Let $p_X(x)$ be the probability mass function of a random variable $X$. Let $Y = g(X)$ be a random variable that is a function of $X$. Then, the expected value of $Y$ is
$$
\mathbb{E}[Y] = \mathbb{E}[g(X)] = \sum_x g(x)p_X(x)
$$
\end{theorem}
\begin{proof}
(Adapted from \href{https://en.wikipedia.org/wiki/Law_of_the_unconscious_statistician}{Wikipedia}) Assume that $X$ takes on finitely many different values $x_1, x_2, ...$. The random variable $g(X)$ then has values $g(x_1), g(x_2), ...$, although some of these may coincide with each other. Let $y_1, y_2, ...$ enumerate the possible distinct values of $g(X)$, and for each $i$ let $I_i$ denote the collection of all $j$ with $g(x_j) = y_i$. Let $p_Y(y)$ denote the probability mass function of $Y$. Since a $y_i$ can be the image of multiple distinct $x_j$,
$$
p_Y(y_i) = \sum_{j \in I_i} p_X(x_j)
$$
Then, according to the definition of expected value,
$$
\begin{aligned}
\mathbb{E}[Y] = \mathbb{E}[g(X)]
&= \sum_i y_i p_Y(y_i) \\
&= \sum_i y_i \sum_{j \in I_i} p_X(x_j) \\
&= \sum_i \sum_{j \in I_i} y_i p_X(x_j) \\
&= \sum_i \sum_{j \in I_i} g(x_j) p_X(x_j) \\
&= \sum_x g(x) p_X(x)
\end{aligned}
$$

Following the conventions of CS 229 and Wikipedia, the notation used to explicitly indicate that the expected value of $Y = g(X)$ is expressed in terms of the probabiltiy distribution of $X$ is $\mathbb{E}_{X \sim p_X}[g(X)]$.
\end{proof}

\begin{definition}[Evidence lower bound (ELBO)]
Let $X$ and $Z$ be random variables jointly distributed with distribution $p(x, z \mid \theta)$. Then, for any distribution $q(z \mid x, \phi)$
$$
\begin{aligned}
\textbf{ELBO}(x, q, \theta)
&= \sum_z q(z \mid x, \phi) \log \frac{p(x, z \mid \theta)}{q(z \mid x, \phi)} \\
&= \mathbb{E}_{z \sim q(z \mid x, \phi)} \left[\log \frac{p(x, z \mid \theta)}{q(z \mid x, \phi)} \right]
\end{aligned}
$$
where the second equality relies on the law of the unconscious statistician.
\end{definition}

\begin{theorem}[Jensen's inequality]
Let $f$ be a convex function, and let $X$ be a random variable. Then $\mathbb{E}[f(X)] \geq f(\mathbb{E}[X])$. Moreover, if $f$ is strictly convex\footnote{A strictly convex scalar function is defined by $f''(x) > 0$ for all $x$. A strictly convex vector-valued function has a positive definite Hessian, $H > 0$.}, then $\mathbb{E}[f(X)] = f(\mathbb{E}(X))$ if and only if $X$ is a constant ($X = \mathbb{E}[X]$).
\end{theorem}

As a concrete example, if $f$ is a scalar convex function, and $X$ is a scalar random variable with probability mass function $p_X(x)$, Jensen's inequality states that
$$
\sum_x \left(f(x) \cdot p_X(x)\right) \geq f\left(\sum_x \left( x \cdot p_X(x) \right) \right).
$$
Applying the law of the unconscious statistician gives the shorthand notation $\mathbb{E}[f(X)] \geq f(\mathbb{E}[X])$.

\subsection{Problem setup}

Given the statistical model which generates a set $\mathbf{X}$ of observed data, a set of unobserved latent data or missing values $\mathbf{Z}$, and unknown parameters $\theta$ (may be a vector), along with a likelihood function
$$
L(\theta; \mathbf{X}, \mathbf{Z}) = p(\mathbf{X}, \mathbf{Z} \mid \theta)
$$
the maximum likelihood estimate (MLE) of the unknown parameters is
$$
\argmax_\theta L(\theta; \mathbf{X})
= \argmax_\theta \sum_\mathbf{Z} p(\mathbf{X}, \mathbf{Z} \mid \theta)
= \argmax_\theta \sum_\mathbf{Z} p(\mathbf{X} \mid \mathbf{Z}, \theta) p(\mathbf{Z} \mid \theta)
$$

\section{CS 229}

Consider the log-likelihood for a single sample $x$ (and its corresponding latent variable $z$). Let $Q$ be a \textbf{proposal distribution} over possible values of $z$.
$$
\begin{aligned}
\log p(x \mid \theta)
&= \log \sum_z p(x, z \mid \theta) \\
&= \log \sum_z Q(z) \frac{p(x, z \mid \theta)}{Q(z)} \\
&\geq \sum_z Q(z) \log \frac{p(x, z \mid \theta)}{Q(z)} \\
&= \text{ELBO}(x, Q, \theta)
\end{aligned}
$$

where the last step used Jensen's inequality. Specifically, $f(x) = \log x$ is a concave function ($f''(x) = - 1/x^2 < 0$ over its domain $x \in \mathbb{R}_+$). Then,

$$
f\left(\mathbb{E}_{z \sim Q} \left[\frac{p(x, z \mid \theta)}{Q(z)}\right]\right)
\geq \mathbb{E}_{z \sim Q} \left[f\left( \frac{p(x, z \mid \theta)}{Q(z)}\right) \right]
$$

where the $z \sim Q$ subscripts indicate that the expectations are with respect to $z$ drawn from $Q$. To be more explicit about how Jensen's inequality is used here, let $g(Z)$ be the random variable
$$
g(Z) = \frac{p(x, Z \mid \theta)}{Q(Z)}.
$$
Note that $Z$ is drawn according to $Q$. Then, Jensen's inequality is applied in the form
$$
f(\mathbb{E}[g(Z)]) \geq \mathbb{E}[f(g(Z))]
$$

This inequality ($\log p(x \mid \theta) \geq \text{ELBO}(x, Q, \theta)$) holds for any distribution $Q$. We can make this inequality tight (i.e., hold with equality) if (as noted in Jensen's inequality theorem) $g(Z)$ is constant: let $Q(z) \propto p(x, z \mid \theta)$. Since $Q$ is a distribution and $\sum_z Q(z) = 1$,
$$
\begin{aligned}
Q(z)
&= \frac{p(x, z \mid \theta)}{\sum_z p(x, z \mid \theta)} \\
&= \frac{p(x, z \mid \theta)}{p(x \mid \theta)} \\
&= p(z \mid x, \theta)
\end{aligned}
$$

The EM algorithm alternatively updates $Q$ and $\theta$:
\begin{itemize}
    \item E-step: Set $Q$ to be the posterior distribution of $z$ given $x$ and the \textit{current} parameters.
    $$
    Q(z) = p(z \mid x, \theta^{(t)})
    $$
    To be explicit about the dependencies of $Q$, we can write $Q(z \mid x, \theta^{(t)})$ instead of $Q(z)$. This matches the form of $q(z \mid x, \phi)$ used in the definition above of the ELBO, with $\phi = \theta^{(t)}$.
    \item M-step: Maximize the ELBO function with respect to $\theta$ while \textit{fixing} the choice of $Q$.
    $$
    \begin{aligned}
    \theta^{(t+1)}
    &= \argmax_\theta \text{ELBO}(x, Q, \theta) \\
    &= \argmax_\theta \sum_z Q(z \mid x, \theta^{(t)}) \log \frac{p(x, z \mid \theta)}{Q(z \mid x, \theta^{(t)})} \\
    &= \argmax_\theta \sum_z Q(z \mid x, \theta^{(t)}) \log p(x, z \mid \theta) \underbrace{- \sum_z Q(z \mid x, \theta^{(t)}) \log Q(z \mid x, \theta^{(t)})}_{\text{entropy of } Z: \, H(Z \mid x, \theta^{(t)})} \\
    &= \argmax_\theta \sum_z Q(z \mid x, \theta^{(t)}) \log p(x, z \mid \theta) \\
    &= \argmax_\theta \mathbb{E}_{z \sim Q} \log p(x, z \mid \theta)
    \end{aligned}
    $$
    \begin{itemize}
        \item Notice how the entropy of $Z \mid x, \theta^{(t)}$ does not depend on $\theta$, so that term can be ignored. This makes the CS 229 expression for the M-step equivalent to the Wikipedia expression for the M-step.
        \item The M-step can be interpreted as maximizing the expected value (over the proposal distribution) of the ``complete-data'' log-likelihood.
    \end{itemize}
    \item CS 229 appears to follow an unconventional split of the E- and M-steps. In the conventional description of the EM algorithm (both in the orignal \href{https://doi.org/10.1111/j.2517-6161.1977.tb01600.x}{Dempster et al. (1977) paper} and in Wikipedia), the E-step (\textit{expectation step}) defines the expected value function
    $$
    \mathbb{Q}(\theta \mid \theta^{(t)}) = \mathbb{E}_{z \sim Q} \log p(x, z \mid \theta) = \sum_{z} Q(z \mid x, \theta^{(t)}) \log p(x, z \mid \theta)
    $$
and the M-step (\textit{maximization step}) finds the parameters $\theta$ that maximize that expected value function
$$
\theta^{(t+1)} = \argmax_\theta \mathbb{Q}(\theta \mid \theta^{(t)}).
$$
\end{itemize}

Generalized to a dataset consisting of multiple samples $\{x^{(1)}, ..., x^{(n)}\}$, the log likelihood becomes
$$
\begin{aligned}
\ell(\theta)
&= \sum_{i=1}^n \log p(x^{(i)} \mid \theta) \\
&= \sum_{i=1}^n \log \sum_{z^{(i)}} p(x^{(i)}, z^{(i)} \mid \theta) \\
&\geq \sum_{i=1}^n \sum_{z^{(i)}} Q_i(z^{(i)}) \log \frac{p(x^{(i)}, z^{(i)} \mid \theta)}{Q_i(z^{(i)})} \\
&= \sum_{i=1}^n \text{ELBO}(x^{(i)}, Q_i, \theta)
\end{aligned}
$$
Here, we introduce one distribution $Q_i$ for each samples $x^{(i)}$ and set each $Q_i$ to be the posterior distribution
$$
Q_i(z^{(i)}) = p(z^{(i)} \mid x^{(i)}, \theta)
$$

The full EM algorithm is therefore as follows:
\begin{enumerate}
    \item Initialize $\theta^{(0)}$.
    \item Repeat until convergence (either the ELBO stops improving, or the parameter values stop changing):
    \begin{itemize}
        \item E-step: For each $i$, set 
        $$
        Q_i(z^{(i)}) = p(z^{(i)} \mid x^{(i)}, \theta^{(t)})
        $$
        \item M-step: Set
        $$
        \begin{aligned}
        \theta^{(t+1)}
        &= \argmax_\theta \sum_{i=1}^n \text{ELBO}(x^{(i)}, Q_i, \theta) \\
        &= \argmax_\theta \sum_{i=1}^n \sum_{z^{(i)}} Q(z^{(i)} \mid x^{(i)}, \theta^{(t)}) \log \frac{p(x^{(i)}, z^{(i)} \mid \theta)}{Q(z^{(i)} \mid x^{(i)}, \theta^{(t)})} \\
        &= \argmax_\theta \sum_{i=1}^n \sum_{z^{(i)}} Q(z^{(i)} \mid x^{(i)}, \theta^{(t)}) \log p(x^{(i)}, z^{(i)} \mid \theta) \\
        &= \argmax_\theta \sum_{i=1}^n \mathbb{E}_{z^{(i)} \sim Q_i} \log p(x^{(i)}, z^{(i)} \mid \theta)
        \end{aligned}
        $$
    \end{itemize}
\end{enumerate}

\section{BE/Bi 103b}

While CS 229 describes EM primarily as a substitute for MLE when there are latent variables, it can be used to find other estimators, such as the maximum a posteriori (MAP) estimator. This is the approach taken in the BE/Bi 103b notes and is also briefly mentioned in the original \href{https://doi.org/10.1111/j.2517-6161.1977.tb01600.x}{Dempster et al. (1977) paper}.

The marginal posterior distribution of interest is
$$
p(\theta \mid x)
= \frac{p(x, \theta)}{p(x)}
$$

The MAP estimate is the value of $\theta$ that maximizes $p(\theta \mid x)$. Since the denominator $p(x)$ is not a function of $\theta$, the objective function is
$$
\begin{aligned}
h(\theta)
&= \log p(x, \theta) \\
&= \log \sum_z p(x, z, \theta) \\
&= \log \sum_z p(x, z \mid \theta) p(\theta) \\
&= \log \sum_z Q(z) \frac{p(x, z \mid \theta) p(\theta)}{Q(z)} \\
&\geq \sum_z Q(z) \log \frac{p(x, z \mid \theta) p(\theta)}{Q(z)} \\
\end{aligned}
$$
where $p(\theta)$ is the prior distribution of the parameter values. To make this inequality tight, set $Q(z) = p(z \mid x, \theta)$ such that
$$
\begin{aligned}
\sum_z Q(z) \log \frac{p(x, z \mid \theta) p(\theta)}{Q(z)}
&= \sum_z p(z \mid x, \theta) \log \frac{p(x, z \mid \theta) p(\theta)}{p(z \mid x, \theta)} \\
&= \sum_z p(z \mid x, \theta) \log p(x, \theta) \\
&= \log p(x, \theta) \sum_z p(z \mid x, \theta) \\
&= \log p(x, \theta)
\end{aligned}
$$
where the second equality comes from the fact that $p(x, z, \theta)$ is equal to both $p(z \mid x, \theta) p(x, \theta)$ and $p(x, z \mid \theta) p(\theta)$, so we can express $p(x, \theta)$ as
$$
p(x, \theta) = \frac{p(x, z \mid \theta) p(\theta)}{p(z \mid x, \theta)}.
$$

Thus, the E- and M-steps are
\begin{itemize}
    \item E-step: Set $Q(z \mid x, \theta^{(t)}) = p(z \mid x, \theta^{(t)})$.
    \item M-step: Maximize the surrogate function
    $$
    \begin{aligned}
    \theta^{(t+1)}
    &= \argmax_\theta \sum_z Q(z \mid x, \theta^{(t)}) \log \frac{p(x, z \mid \theta) p(\theta)}{Q(z \mid x, \theta^{(t)})} \\
    &= \argmax_\theta \sum_z Q(z \mid x, \theta^{(t)}) \log p(x, z \mid \theta) p(\theta) \underbrace{- \sum_z Q(z \mid x, \theta^{(t)}) \log Q(z \mid x, \theta^{(t)})}_{\text{entropy of } Z: \, H(Z \mid x, \theta^{(t)})} \\
    &= \argmax_\theta \sum_z Q(z \mid x, \theta^{(t)}) \log p(x, z \mid \theta) p(\theta) \\
    &= \argmax_\theta \sum_z Q(z \mid x, \theta^{(t)}) \left(\log p(x, z \mid \theta) + \log p(\theta) \right) \\
    &= \argmax_\theta \sum_z Q(z \mid x, \theta^{(t)}) \log p(x, z \mid \theta) + \sum_z Q(z \mid x, \theta^{(t)}) \log p(\theta) \\
    &= \argmax_\theta \sum_z Q(z \mid x, \theta^{(t)}) \log p(x, z \mid \theta) + \log p(\theta) \sum_z Q(z \mid x, \theta^{(t)}) \\
    &= \argmax_\theta \log p(\theta) + \sum_z Q(z \mid x, \theta^{(t)}) \log p(x, z \mid \theta) + \\
    &= \argmax_\theta \log p(\theta) + \mathbb{E}_{z \sim Q} \log p(x, z \mid \theta)
    \end{aligned}
    $$
\end{itemize}

\end{document}